\chapter{Proposed Method and Implementation}
\label{ch:method}

This chapter describes the structure of the delivered system, the data
pipeline, the interruption detector, the design of the direct-track experiments,
and the evidence-first way of working.

\section{Early Attempt and Revision}
\label{sec:early-attempt}

The first implementation was based on LLaMA-Omni2: a frozen Whisper encoder, a
projector, the Qwen2.5-0.5B~\cite{qwen2025qwen25} language model with a LoRA
adapter, and separate recognition and generation heads. The loss decreased, but
an internal audit found three fundamental defects: the language model did not
take part in the forward pass; the recognition target used per-process random
hashes; and the generation target was reconstruction of the input utterance,
not a conversational reply. That artifact was not a speech-to-speech model. Two
decisions followed from this audit: serving moved to a cascade with real
components, and for the direct track only a model with an official trainer was
accepted.

\section{Dataset Construction}
\label{sec:data}

No public Persian corpus with interruption labels was available. The dataset was
derived from an internal archive of spoken content from four channels. Internal
research use was recorded with the supervisor's approval; this approval is not a
licence to redistribute the raw audio or captions, and no restricted data are in
the code repository.

In the first design, text labels were produced by re-running the speech
recognition service. That same service had earlier been tuned on the captions of
these channels, which created a circular loop. Once this was identified, the
text labels were switched to the original source captions, and recognition was
kept only as a runtime component and a measurement tool.

The pipeline is: a caption inventory of 775.887 hours, weighted programme
selection, reconstruction of fifteen-minute windows, speaker diarization, noise
condition labels, extraction of ``user speech / response'' pairs without reuse,
and an audited stereo export with channel zero for the reconstructed assistant
and channel one for the user. Train/validation/test splits are by source session
identifier. The interaction labels are automatic candidates, and the number of
human-verified rows is zero.

\section{Architecture of the Streaming Cascade}
\label{sec:cascade}

The delivered system is a streaming cascade with full-duplex control
(Figure~\ref{fig:arch}). The first version worked with a push-to-talk button and
built the whole reply before playback (about 2.2 seconds to the first audio). In
the final version the stages overlap: the end of speech is detected
automatically, generation starts speculatively, tokens arrive as a stream, and
each clause is synthesized and played as soon as it is complete.

\begin{figure}[ht]
\centering
\scalebox{0.92}{%
\begin{tikzpicture}[
  box/.style={draw,rounded corners=2pt,align=center,font=\scriptsize,
              minimum height=9mm,inner sep=3pt,text width=23mm},
  client/.style={box,fill=black!5},
  lbl/.style={font=\scriptsize,inner sep=1.5pt,fill=white},
  ar/.style={-{Latex[length=1.8mm]},thin},
  dashar/.style={ar,dashed}
]
\node[client] (mic) at (0,0) {browser mic\\20 ms PCM16 frames\\+ playback level};
\node[box] (vad)  at (3.2,-2.3) {Silero VAD endpointer\\speculative 96 ms\\confirm 256 ms};
\node[box] (asr)  at (6.4,-2.3) {NeMo v4 CTC\\ASR};
\node[box] (llm)  at (9.6,-2.3) {Qwen3-4B (vLLM)\\streamed, Persian-only};
\node[box] (tts)  at (12.8,-2.3) {clause chunker\\+ Mana Piper (CPU)};
\node[box] (dtd)  at (3.2,0) {double-talk detector\\+ Silero: \emph{duck}};
\node[box] (lex)  at (6.4,0) {lexical confirm\\(ASR, echo words,\\listener tokens)};
\node[client] (spk) at (12.8,0) {browser playback\\queue: duck / stop};
\draw[ar] (mic) |- (vad.west) node[lbl,pos=0.75] {listening};
\draw[ar] (vad) -- (asr);
\draw[ar] (asr) -- (llm);
\draw[ar] (llm) -- (tts);
\draw[ar] (tts) -- (spk) node[lbl,midway] {per clause};
\draw[ar] (mic) -- (dtd) node[lbl,midway,above=1pt] {while speaking};
\draw[ar] (dtd) -- (lex);
\draw[ar] (lex) -- (spk) node[lbl,midway] {stop / resume};
\draw[dashar] (lex.south) -- ++(0,-0.6) -| (vad.north east)
      node[lbl,pos=0.3] {stop: 200 ms pre-roll + speech to next turn};
\end{tikzpicture}}
\caption{The streaming architecture. The bottom row is reply generation while
         listening; the top row is interruption control while the assistant
         speaks.}
\label{fig:arch}
\end{figure}

The browser client captures the microphone continuously and sends one frame
every 20 ms (a sequence number, the level of the assistant audio being played,
and 320 PCM16 samples); it schedules the reply clauses on a playback queue and
reports timings with its own clock. The Silero detector starts a speculative turn
after 96 ms of silence and holds its audio until the end of speech is confirmed
at 256 ms; if the user continues, the turn is cancelled and the speech is added
to the same turn. Recognition uses NeMo version 4 (CTC), which the candidate
comparison found to be the most accurate candidate within the latency budget.
The reply is generated by Qwen3-4B-Instruct-2507~\cite{qwen2025qwen3} with the
locked v2 prompt, streamed from vLLM with two turns of history, and a
``Persian script only'' constraint is applied during decoding through
structured output. The text is cut at clause boundaries (a first fragment of 2--4
words, then at most 14 words), and each fragment is synthesized with Piper and
the Mana voice on the CPU.

The nine-row mechanical panel and the earlier re-recognition measure were run on
the chain before the upgrade (Qwen2.5-0.5B-Instruct); the 40-row semantic test
is on Qwen3-4B with the Qwen3.8-27B-FP8 judge.

\section{Interruption Detection and the Full-Duplex Layer}
\label{sec:bargein}

While the assistant speaks, a two-stage controller inspects the frames. Stage
one: a frame must be speech and louder, by a margin, than the expected echo of
the assistant's voice (the playback level reported by the client plus an
estimated coupling factor, assumed conservatively at the start of each reply and
then adapted); one such frame ducks the assistant, giving a ``detect and adapt''
reaction with a 150 ms target. Stage two: the speech captured since the onset of
the interruption is transcribed by the same recognizer; listener tokens
(\fa{آره}, \fa{بله}, \fa{اوهوم}, \ldots), words of the reply being played
(residual echo) and function words are discarded. If any word remains, playback
stops, generation is cancelled, and the user's speech, with a 200 ms pre-roll,
becomes the next turn; otherwise the volume returns and the reply continues.
Development runs of the simulation showed that a confirmation based on speech
duration or on an acoustic classifier does not separate long backchannels, loud
noise and echo from interruptions.

Evaluation has three levels: a classifier test on podcast events with
independent two-pass automatic Gemini labels, a deployment simulation built from
real recordings, and loopback testing with a real browser. The test
backchannels were selected with a different recognition model
(SeamlessM4T~v2), so that the lexical confirmation is not circular.

\section{Design of the Direct-Track Experiments}
\label{sec:direct}

The direct track is low-rank adaptation of Moshika with the official trainer. The
input is the stereo export of Section~\ref{sec:data}. The loss is a weighted
combination of the text loss and the audio-token loss.

The runtime acceptance condition was locked before training: the adapter is
loaded in the official streaming server, and for each of nine predefined rows
there must be text, non-silent speech and the Persian-script constraint. A
checkpoint is eligible only with 9/9. The condition is deliberately strict,
because the lowest validation loss in the first version gave almost silent audio
without text.

The experiments are a chain of falsifiable hypotheses, and in each generation
mostly one factor changed (Table~\ref{tab:moshi-design}). Version v6.2 is a
capacity test, not a generalization test: training and evaluation on the same
32 samples, a 12-second window, and scheduled dropout of the text input from
0.25 to 0.75. If the loss does not decrease even on training data, the problem
is capacity; if it decreases but runtime fails, the problem is exposure bias.

\begin{table}[ht]
\centering
\caption{Design of the seven experimental generations of the direct track}
\label{tab:moshi-design}
\small
\begin{tabular}{lL{3.6cm}L{5.0cm}c}
\toprule
\textbf{Version} & \textbf{Changed factor} & \textbf{Tested hypothesis} &
\textbf{Rank} \\
\midrule
v1 & base configuration, 20-second window &
the model learns Persian speech with enough data & 64 \\
v2 & 100-second text window &
a short window does not fit the start of the reply in context & 64 \\
v3 & frozen embeddings, higher rank &
training the embeddings damages the base representation & 128 \\
v4 & only the two text embeddings &
only the Persian text space needs adaptation & 128 \\
v5 & first-codebook loss weight from 100 to 10 &
a heavy semantic-codebook weight stifles acoustic learning & 128 \\
v6.2 & in-sample capacity test, 12-second window &
the model has capacity and the problem is exposure bias & 64 \\
v7a & romanized phoneme text, word-level timing &
the tokenizer breaks the one-token-per-frame budget & 64 \\
\bottomrule
\end{tabular}
\end{table}

The final audit revealed a structural cause that the first six generations had
not tested: Moshi's English tokenizer needs 6.6 tokens per Persian word (55\%
byte fallback), while the model produces at most one text token per 12.5 Hz
frame; 98.5\% of turns exceeded this budget. Version v7a made the inner-monologue
text a romanization of the phonemes Piper actually spoke (2.7 tokens per word)
with word-by-word timing, and kept the other conditions of the capacity test
fixed. The second direct track is Gemma~4~E4B, which hears audio directly; it was
measured against the cascade zero-shot and with a DiVA-style self-distillation
LoRA adapter (3,817 turns, 11.2 hours).

\clearpage
\section{Evidence-First Method}
\label{sec:evidence-first}

Five operating rules were applied. One: the protocol, row indices and thresholds
are registered before running. Two: semantic evaluation has two stages; only
after the seven-row qualification passes completely is the final test opened,
once. The Qwen3.8-27B-FP8 judge runs at temperature zero, in blind order and
with two repetitions. Three: result artifacts are write-once. Four: each
measurement has an evidence class (Table~\ref{tab:evidence-classes}), and only
the official class can prove a requirement of the project definition. Five: the
aggregator fails closed; a missing artifact or a class below what is needed makes
a condition fail, not ``unknown''.

\begin{table}[ht]
\centering
\caption{Evidence classes and the claim power of each class}
\label{tab:evidence-classes}
\small
\begin{tabular}{L{3.6cm}L{6.0cm}c}
\toprule
\textbf{Evidence class} & \textbf{Required conditions} & \textbf{Proves a requirement} \\
\midrule
Official end-to-end & live consented session, client playback
confirmation, real microphone and eligible hardware & yes \\
Unofficial live & same protocol, without eligible hardware or study conditions &
no \\
Server-side component & generation or detector timing, without proof of the
client's audible behaviour & no \\
Synthetic approximation & generated audio or simulated events & regression test only \\
Same-family automatic approximation & judging by an LLM of the candidate's
family & no \\
\bottomrule
\end{tabular}

\vspace{0.5em}
{\footnotesize\itshape The consequence of this method is that the evidence
status of each condition is stated independently of whether the report is ready
for presentation and examination. The final result is a clear separation of what
was built from what was proven with official evidence.\par}
\end{table}
